\documentclass[12pt]{article} % specifies font size and document type
\usepackage{amsmath,amsthm,amssymb,amsfonts,mathtools} % these packages are necessary for math formatting
\usepackage{geometry}
\usepackage[scheme=plain]{ctex}
\usepackage{tikz-cd}

\geometry{tmargin=.75in, bmargin=.75in, lmargin=.75in, rmargin = .75in}  

\theoremstyle{definition}
\newtheorem{defn}{Definition}[subsection]
\newtheorem{prop}{Proposition}[subsection]
\newtheorem{theorem}{Theorem}[subsection]
\newtheorem{rem}{Remark}[subsection]
\newtheorem{lem}{Lemma}[subsection]
\newtheorem{cor}{Corollary}[subsection]

\newcommand{\R}{\mathbb{R}}
\newcommand{\C}{\mathbb{C}}
\newcommand{\Z}{\mathbb{Z}}
\newcommand{\N}{\mathbb{N}}
\newcommand{\Q}{\mathbb{Q}}
\newcommand{\der}{\mathrm{d}}
\newcommand{\vr}{\vec{r}}


\begin{document}

\title{Notes for Differential Geometry}
\author{S. Y.H.}
\date{July 3, 2026}

\maketitle

\newpage
\tableofcontents

%%%%%%%%%%%%%%%%%%%%%%%%%%%%%%%%%%%%%%%%
\newpage
\section{Local theory of curves} 

\begin{defn}
    A curve \(\vr=\vr(t)\colon (a,b)\rightarrow\R^n\) where \(\vr(t)=(x_1(t),\cdots,x_n(t))\) is \textbf{regular} if 
    \(x_i(t)\in C^{\infty}(a,b)\) and for all \(t\in(a,b)\), we have \(\left|\vr\,'(t)\right|\neq 0\).
\end{defn}

Specifically, we often take \(n=2,3\).

We introduce the following notations. For a general variable like \(t\), \(\vr\,'=\frac{\der}{\der t}\vr(t)\).

\begin{defn}
    The arc length of \(\vr\), denoted by \(s\), is defined by
    \[s=\int_{a}^{t}{\left|\vr\,'(\tau)\right|\der\tau}.\]
\end{defn}

Hence, we have \(\frac{\der s}{\der t}=\left|\vr\,'(t)\right|.\)
We use a different notation to denote the derivative of \(\vr\) with respect to \(s\): \(\dot{\vr}(s)=\frac{\der}{\der s}\vr(s)\).
The benefit of using the arc length as the variable is the normalization of the tangent vector. That is,
\[\left|\dot{\vr}(s)\right|=\dfrac{\der \vr}{\der s}=\dfrac{\der \vr}{\der t}\dfrac{\der t}{\der s}=\vr\,'\cdot\dfrac{1}{\left|\vr\,'\right|}=1.\]

Let \(T(s)=\dot{\vr}(s)\) be the tangent vector and \(N(s)\) be the normal vector with \(\left|N\right|=1\).
We always assume that \(T\) is rotated counterclockwise by \(\pi/2\) to obtain \(N\).

\subsection{Plane curves}
Let \(\vr=\vr(s)\), \(T(s)=\dot{\vr}(s)\) with \(\left|T(s)\right|=1\).
Then \(\left|T(s)\right|^2=1\). We have \(\frac{\der}{\der s}\left|T(s)\right|^2=0\). i.e. \(\frac{\der}{\der s}\langle T(s),T(s)\rangle=2\langle T,\dot{T}\rangle=0\).
\newpage
\section{Local theory of surfaces}
    

\newpage
\section{Exterior differential systems}

\newpage



\end{document}